\section*{集合论的悖论与基础危机。}
\phantomsection
\addcontentsline{toc}{section}{集合论的悖论与基础危机。}

最初的“悖论的”集合出现在基数和序数的理论中。1897 年，布拉利-福尔蒂指出，不能考虑由全部序数组成的集合的存在，因为这个集合将是良序的，从而同构于它的一个异于全体的段，而这是荒谬的 {[}42{]}。\(^{55}\) 1899 年，康托尔（在一封致戴德金的信中）注意到，无论谈论基数组成一个集合，还是谈论“所有集合的集合”，都不再可能不陷入矛盾（这最后一个“集合” \(\Omega\) 的子集的集合将与 \(\Omega\) 的一个子集等势，这与不等式 \(m < 2^{m}\) 相悖）（{[}47{]}，第~444-448~页）。最后，1905 年，罗素在分析这个不等式的证明时表明，建立该不等式的推理（无须求助于基数论）同时也证明了“所有不是其自身元素的集合所组成的集合”这一概念是矛盾的 {[}266{]}。\(^{56}\)

人们或许会以为，这样的“二律背反”只会出现在数学的边缘地带，即出现在以考虑直觉无法企及的“大小”的集合为特征的区域。然而另一些“悖论”很快就威胁到数学中最经典的部分。贝里和罗素（{[}266{]}，第 I 卷，第~63-64~页）简化了 J. 里夏尔的一个论证 {[}258{]}，他们实际上注意到：定义可用英语少于十六个词表述的整数所组成的集合是有限的，然而把一个整数定义为“用英语少于十六个词所不能定义的最小整数”却是矛盾的，因为这个定义本身只由十五个词组成。

这类推理与数学家通常使用的推理相距甚远，在许多人看来也许只像一种字谜，但它们毕竟指出了修改数学基础、以期消除这类“悖论”的必要性。然而，尽管关于这种修改的紧迫性人人意见一致，关于实现它的方式，很快就出现了根本的分歧。对第一派数学家——无论他们是“观念论者”还是“形式主义者”\(^{57}\)——来说，集合论“悖论”所造成的局面，与几何学中因发现非欧几何或“病态”曲线（如无切线的曲线）而产生的局面十分类似；由此应当得出一个类似但更为一般的结论，即想把任意一门数学理论建立在（明确的或不明确的）对“直觉”的诉诸之上，是毫无希望的。这一立场可以用形式主义学派主要论敌的原话来概括：“形式主义者”，布劳威尔说（{[}39 a{]}，第~83~页），“认为人的理性并不拥有例如直线或大于十的数的精确图像\ldots 诚然，对于数学实体之间的某些关系，即我们当作公理的那些关系，我们可以按照固定的规则从其他关系把它们推导出来，并相信这样一来我们便是通过逻辑推理从其他真理导出了真理\ldots{[}但是{]}对形式主义者来说，数学的精确性只存在于关系串的发展过程之中，它独立于人们可能愿意赋予这些关系或它们所联结的实体的任何意义。”

因此，对形式主义者来说，事情就在于给集合论一个与初等几何十分类似的公理基础：在那里，人们不必关心被称为“集合”的那些“对象”是什么，也不必关心关系 \(x \in y\) 意味着什么，而只需列举出应加于这个关系上的条件；当然，这样做时必须尽可能把康托尔理论的全部结果都包括进来，同时又使“悖论的”集合不可能存在。这种公理化的第一个例子由策梅洛于 1908 年给出 {[}342 c{]}；他通过引入“分离公理（Aussonderung）”来避开“过大的”集合：性质 \(P(x)\) 仅当 \(P(x)\) 已蕴涵某个形如 \(x \in A\) 的关系时，才确定出一个由具有该性质的元素组成的集合。⁵⁸ 但是，要消除类似于“里夏尔悖论”的悖论，只能靠限制附加在“性质”概念上的意义来实现；在这一点上，策梅洛只满足于用十分含糊的措辞描述一类他称之为“已定义的”性质，并指出在应用分离公理时必须限于这些性质。这一点由斯科朗 {[}286 a{]} 和弗兰克尔 {[}113 b{]} 加以精确化；正如他们所指出的，对这一点的澄清迫使我们置身于一个完全形式化的系统之中，在那里，“性质”与“关系”这些概念已失去全部“意义”，只是按照明确的规则构造出来的汇集的简单称谓。这当然就要求把所使用的逻辑规则也放进系统之内，而这在策梅洛-弗兰克尔系统中尚付阙如。

在此之后，人们又提出了集合论的其他一些公理化。其中主要应提到冯·诺伊曼的公理化 {[}342 a and c{]}，它比策梅洛-弗兰克尔系统更接近康托尔的原初构想：后者为避开悖论的集合，早在与戴德金的通信中就已提议（{[}47{]}，第~445-448~页）区分两类集合，即“杂多”（“Vielheiten”）与本来意义上的“集合”（“Mengen”），后者的特征在于它们可以被设想为单一的对象。冯·诺伊曼正是把这一想法精确化，区分出两类对象：“集合”与“类”；在他的系统中（几乎是完全形式化的），类与集合的区别在于：类不能被置于记号 ∈ 的左边。这种系统的好处之一，是它为 19 世纪逻辑学家所使用的“全类”概念恢复了名誉（全类自然不是集合）；还应指出，冯·诺伊曼的系统（就集合论而言）避免了公理模式的引入，而代之以合适的公理（这使得对它作逻辑研究更为容易）。贝尔奈斯和哥德尔给出了冯·诺伊曼系统的一些变体 {[}130 b{]}。

上述系统看来已经妥善地消除了悖论，但代价是引入了一些只能让人觉得十分任意的限制。为策梅洛-弗兰克尔系统说句公道话：它所作的限制只不过认可了把集合概念应用于各门数学理论时的现行实践。冯·诺伊曼和哥德尔的系统离通常的观念更远；但另一方面，某些尚处于初创阶段的数学理论，要安插在这样的系统所提供的框架中，说不定反而比安插在策梅洛-弗兰克尔系统较窄的框架中更容易，这并非不可能。

当然，谁也不能断言这些解决方案中的任何一个给人以定论的印象。它们之所以让形式主义者满意，是因为后者拒绝考虑每个数学家的个人心理反应；他们假定，当一种形式化语言能够把数学推理转录为一种没有歧义、从而适于充当数学思想之载体的形式时，它便完成了自己的任务；他们会说，只要一个人的推理能够转录进共同的语言，那么他对数学对象的“本性”或他所使用的定理的“真理性”爱怎么想就怎么想。\(^{59}\)

换句话说，从哲学的观点看，形式主义者的态度就在于对“悖论”所提出的问题不闻不问，从而抛弃了那种力图赋予数学概念以全体数学家共同的理智“内容”的柏拉图式立场。面对这一与传统的决裂，许多数学家望而却步。例如罗素就试图通过对悖论的结构作更深入的分析来避免它们。罗素和怀特海重新拾起由 J. 里夏尔最先提出（在他展示其“悖论”的那篇文章 {[}258{]} 中）、后又由 H. 庞加莱加以发展 {[}251 c{]} 的一个想法，他们注意到，悖论集合的定义全都违反如下这条称为“恶性循环原则”的原则：“其定义涉及一个集合的元素之全体的那种元素，不能属于该集合”（{[}266{]}，第 I 卷，第~40~页）。也正是这一陈述充当了《数学原理》（Principia）的基础，该书正是为了满足它才发展出“类型论”。与启发它的弗雷格的逻辑一样，罗素和怀特海的逻辑含有“命题变元”；类型论对这些不同的变元进行分类，其主要轮廓如下。从一个未加确定的“个体域”出发——其中的个体可称为“0 阶对象”——，凡变元（自由的或约束的）皆为个体的那些关系，称为“一阶对象”；一般地，凡变元为 ≤ n 阶对象（其中至少一个为 n 阶）的那些关系，称为“\(n + 1\) 阶对象”。\(^{60}\) 于是，n 阶对象组成的一个集合便只能由一个 \(n + 1\) 阶的关系来定义，这一条件使悖论集合得以毫无困难地被消除。\(^{61}\) 但“类型层级”的原则限制性极强，若严格坚持它，得到的就是一种复杂得无法梳理的数学。\(^{62}\) 为了摆脱这一后果，罗素和怀特海不得不引入一条“可化归性公理”，它断言：对于“个体”之间的每一个关系，都存在一个与之等价的一阶关系；这条公理与形式主义者的公理同样任意，而且大大削弱了《数学原理》这一构造的意义。因此，罗素和怀特海的系统在逻辑学家那里比在数学家那里更成功；无论如何，它并未完全形式化，\(^{63}\) 因而存在许多晦暗不明之处。人们曾作出各种努力来简化和澄清这一系统（拉姆齐、赫维斯泰克、蒯因、罗塞尔）；这些作者倾向于使用越来越完全形式化的语言，他们用一些只考虑所论聚合体的措辞的限制，取代了《数学原理》的那些规则（后者尚有某种直观的基础）；这些规则不仅看起来与策梅洛-弗兰克尔系统或冯·诺伊曼系统中所表述的禁令同样毫无根据，而且由于离数学家的实践更远，它们在许多情形下导致了其创立者未曾预见的不可接受的后果（例如布拉利-福尔蒂悖论，或对选择公理的否定）。

对前述各派的数学家来说，首要之点在于不放弃过去遗产的任何部分：“谁也不能把我们从康托尔为我们创造的乐园中驱逐出去”，希尔伯特说（{[}163 c{]}，第~274~页）。为此，他们愿意接受对数学推理的限制；这些限制由于与习惯相合而无伤大体，但它们看来并非由我们的心智习惯和集合概念的直观观念所强加。在他们看来，任何事情都好过让心理学混入数学有效性的判据；与其像阿达马所说的那样“把我们大脑的特性考虑在内”（{[}12{]}，第~270~页），他们宁可听任给数学领域强加一些大体上任意的界限，只要这些界限把经典数学包括在内，又不至于妨碍未来的进步。

而当我们要看的是我们行将谈及的那一倾向的数学家的态度时，情况就完全不同了。如果形式主义者甘愿放弃对数学推理的“心灵之眼”的监督，那么被称为“经验主义者”、“实在论者”或“直觉主义者”的数学家则拒绝这一弃权；他们必须有一种内在的确定性，为保证他们所处理的数学对象的“存在”作保。当问题只在于放弃空间直观时，还没有什么严重的异议，因为算术“模型”容许人们退守到整数这个直观概念后面。但当要把整数概念化归为（直观上远不那么清晰的）集合概念、进而给集合的操作加上毫无直观基础的屏障时，不可调和的对立就显露出来了。这些反对者中最早的一位（而且凭其天才的权威，也是影响最大的一位）是 H. 庞加莱；他不仅承认几何学的公理学观点和分析的算术化，还承认康托尔理论的相当一部分（他是最早在其工作中富有成效地应用这一理论的人之一），但他却拒绝设想算术本身也能通过公理化的处理得到论证；在他看来，完全归纳原理尤其是我们心智的一种基本直觉，不可能把它看成一种纯粹的约定\(^{64}\) {[}251 c{]}。他从原则上敌视形式化语言，并质疑其用处；他不断地把形式化数学中的整数概念，与当时刚刚兴起、我们将于后文谈到的证明论中对整数的使用混为一谈；毫无疑问，在那个时代，要像今天这样——经过 50 年的研究和讨论之后——把这一区别弄清楚是很困难的，尽管一个希尔伯特或一个罗素当时已清楚地感到了它。

这类批评在策梅洛于 1904 年引入选择公理 {[}342 a{]} 之后便成倍增加。在此以前，它在分析或集合论的许多证明中的使用几乎无人觉察；\(^{65}\) 正是遵循埃哈德·施密特提示的一个想法，策梅洛把这条公理明确地表述出来，并以灵巧的方式由它导出了每个集合上都存在良序的一个令人满意的证明。由于与“悖论”同时出现，这种新的推理方式凭其古怪的姿态看来使许多数学家陷入惶惑；只要看看随后一卷《数学年刊》上署名发表的种种奇怪的误解即可，署名者竟是像舍恩弗利斯和 F. 伯恩斯坦那样熟悉康托尔方法的人。E. 博雷尔发表在同一卷上的批评更有分量，而且清楚地附和了 H. 庞加莱关于整数所表达的观点；这些批评在 E. 博雷尔、贝尔、阿达马和勒贝格——他们仍置身于法国数学的经典传统之中——的通信往来中得到展开和讨论 {[}12{]}。博雷尔一开始就否认选择公理的有效性，因为它一般包含不可数无穷多个选择，这在直观上是无法设想的。但阿达马和勒贝格指出，可数无穷多个任意的相继选择也并不更直观，因为它由无穷多次运算组成，而这些运算不可能被设想为实际发生。对把讨论范围扩大了的勒贝格来说，一切都归结为弄清说一个数学对象“存在”是什么意思；在他看来，必须“指名”一个以唯一方式定义该对象的性质（我们会说这是一种“泛函式的”性质）；对于策梅洛在其推理中所使用的那类函数，勒贝格称之为选择的“规律”；他继续说，如果这一要求得不到满足，如果人们只是“想到”这个函数而不“指名”它，那么能肯定在推理过程中人们想到的一直是同一个东西吗（{[}12{]}，第~267~页）？无论如何，这给勒贝格带来了新的疑虑；在集合中选取单个元素的问题在他看来就已引起困难：必须确信这样的元素“存在”，也就是说，能“指名”该集合的至少一个元素。\(^{66}\) 那么，对于一个不知道如何“指名”其每个元素的集合，还能谈论它的“存在”吗？贝尔已经毫不犹豫地否认一个给定无穷集合的所有子集组成的集合的“存在”（出处同上，第~263-264~页）；阿达马指出这些要求将导致放弃谈论实数的可能性，也是徒劳：E. 博雷尔最终赞同的正是这一结论。且不说可数性看来已经取得了永佃权，人们几乎是回到了“实无穷”的反对者的经典立场上。

所有这些异议都谈不上系统；要按类似但更为激进的原则对数学作一次彻底的改造，这任务落到了布劳威尔及其学派的身上。像直觉主义这样一门一半是心理学一半是数学的复杂学说，我们不敢在此加以概述，我们只限于指出其若干最引人注目的方面，欲知详情，请参看布劳威尔本人的著作 {[}39 a and b{]} 以及海廷的阐述 {[}162{]}。对布劳威尔来说，数学与我们思想中的“精确”部分同一，它奠基于自然数序列这一原初直观之上，而且不可能在不遭损毁的情况下把它翻译成一个形式系统。无论如何，它只在数学家的思想中才是“精确”的，想在他们之间发展出一种不受语言的一切缺陷和歧义之累的交流工具，纯属幻想；至多只能希望借助或多或少含糊的描述，在对话者心中唤起一种有利的心态（{[}162{]}，第~11-13~页）。直觉主义数学对逻辑的重视几乎不超过对语言的重视：一个证明之令人信服，并非凭借一劳永逸地固定的逻辑规则，而是凭借其每一步的“直接明证性”。这种“明证性”还必须以比 E. 博雷尔及其追随者更受限制的方式来解释：因此在直觉主义数学中，不能说形如“R 或（非 R）”的关系为真（排中原理），除非对 R 中变元的每一组给定的值，都能证明 R、“非 R”两个命题之一为真；例如，从两个实数之间的方程 ab = 0 不能推出“a = 0 或 b = 0”，因为很容易举出实数 a、b 的明确例子，使得 ab = 0 成立，而在当前却无法证明 a = 0、b = 0 两个命题中的任何一个（{[}162{]}，第~21~页）。

从这样的原则出发，直觉主义数学家最终得到与经典定理大不相同的结果，这不足为怪。这些结果中有整整一部分消失了，例如分析中大部分“存在性的”定理（如实函数的波尔查诺定理和魏尔斯特拉斯定理）；如果一个实变量函数在直觉主义的意义上“存在”，它便理所当然地连续；有界的实数单调序列不一定有极限。另一方面，许多经典概念在直觉主义者那里分岔为若干个不同的基本概念：例如（实数序列的）收敛概念有两个，可数性概念有八个。不言而喻，超限归纳法及其对现代分析的应用（如同康托尔理论的大部分内容一样）被不加申诉地判了死刑。

依照布劳威尔的看法，数学命题只有以这种方式才能获得“内容”；超出直觉主义者容许范围的形式主义推理被判为毫无价值，因为再也无法赋予它们一种可以应用“真”这一直观概念的“意义”了。显然，这样的判断只能依靠一种先已存在的、具有心理学或形而上学本性的“真理”概念。这实际上等于说，它们是拒绝一切讨论的。

毫无疑问，来自直觉主义阵营的猛烈攻击有时不仅迫使数学中的前卫学派、甚至迫使传统数学的追随者采取守势。有一位著名数学家承认，这些攻击给他留下如此深刻的印象，以致他自愿把工作限制在被判定为“可靠”的数学分支之内。但这类情况很难是经常发生的。直觉主义学派的下场，其记忆无疑注定只作为一种历史奇观而留存；但它至少立过一功：它迫使它的论敌——也就是数学家的压倒多数——明确了自己的立场，并更清楚地认识到他们信赖数学的理由（一部分属于逻辑方面，另一部分属于情感方面）。

